\subsection{Routing predittivo e frontiera classico-neurale}
\label{sec:routing-predittivo-frontiera-classico-neurale}

Il router R-D-E non stabilisce una superiorita assoluta dei codec neurali
o dei codec classici. La valutazione operativa descrive invece una
frontiera di switch: per ogni immagine, profilo e vincolo di qualita, il
router confronta il miglior candidato classico e il miglior candidato
neurale nello stesso spazio rate-distortion-energy. I codec neurali
diventano convenienti solo quando il vantaggio di rate o qualita
compensa la penalita energetica. I codec classici restano sufficienti
quando il peso dell'energia, i vincoli percettivi o i vincoli di sistema
dominano la decisione.

\subsubsection{Regimi operativi simulati}

\input{tables/routing_regime/table_regime_summary_core}

La simulazione operativa usa la policy
\texttt{metadata\_plus\_system\_full\_pool} per osservare come la scelta
predittiva cambia al variare del regime. Nel benchmark considerato, il
regime \texttt{bandwidth\_limited} e quello in cui la selezione neurale
rimane piu visibile; nei regimi orientati a energia, pressione termica,
assenza di CUDA o pressione di memoria, il sistema tende invece a
sopprimere o penalizzare i candidati neurali. Questo comportamento non
implica dominanza universale di una famiglia di codec, ma indica che il
profilo operativo modifica in modo sostanziale la convenienza R-D-E.

\subsubsection{Contratto di qualita dell'oracle}

L'oracle minimizza il costo R-D-E solo tra candidati che soddisfano il
quality floor attivo. Se nessun candidato soddisfa il floor, il caso e
marcato come \texttt{infeasible} e non viene trasformato in una scelta
oracle a bassa qualita. Le violation delle policy predittive sono
quindi \emph{realized ex-post violations}: descrivono la qualita
ottenuta dopo che la policy ha scelto senza osservare le misure R-D-E
della target image, e non sono oracle violations.

The oracle is not a low-quality selector: it minimizes the R-D-E
objective only over candidates satisfying the active quality floor, or
reports infeasibility.

\subsubsection{Switch analysis con SSIMULACRA2}

\input{tables/routing_regime/table_ssimulacra2_switch_summary}

Con SSIMULACRA2, i codec neurali vincono solo nel regime
\texttt{bandwidth\_limited} e in una regione ristretta del benchmark.
Al crescere del floor da 60 a 80, \texttt{neural\_win\_rate} passa da
13.54\% a 2.08\%. In tutti questi casi
\texttt{neural\_necessary\_rate} e pari a 0: nel benchmark considerato i
codec neurali non sono necessari per superare la soglia percettiva, ma
sono R-D-E efficienti in una frazione limitata dei casi.

\begin{figure}[t]
  \centering
  \includegraphics[width=0.85\textwidth]{figures/routing_regime/switch_reason_by_rate_weight.png}
  \caption{Distribution of neural/classical switch reasons as the rate weight increases.}
  \label{fig:switch-reason-rate-weight}
\end{figure}

\begin{figure}[t]
  \centering
  \includegraphics[width=0.85\textwidth]{figures/routing_regime/neural_vs_classic_tradeoff_scatter.png}
  \caption{Neural versus classical trade-off between bitrate reduction and energy penalty.}
  \label{fig:neural-classic-tradeoff}
\end{figure}

\begin{figure}[t]
  \centering
  \includegraphics[width=0.85\textwidth]{figures/routing_regime/quality_floor_vs_neural_necessity.png}
  \caption{Effect of the SSIMULACRA2 quality floor on neural usefulness and necessity.}
  \label{fig:quality-floor-neural-necessity}
\end{figure}

\subsubsection{Expected-quality gate}

\input{tables/routing_regime/table_expected_quality_gate_summary}

Il gate shadow elimina le violation ex-post nel run SSIMULACRA2, ma al
costo di una policy piu conservativa: aumenta il fallback, sopprime le
selezioni neurali e puo aumentare il regret medio. Per questo motivo non
viene presentato come policy finale, ma come diagnostica di
safety/adaptivity: mostra come ridurre le violation senza leakage,
usando solo statistiche di qualita stimate dal training fold.

\subsubsection{PSNR30 come stress test}

PSNR30 viene usato solo come stress test rate-oriented. Non e il claim
finale di qualita percettiva e non sostituisce la valutazione basata su
SSIMULACRA2. Nel regime \texttt{bandwidth\_limited}, PSNR30 mostra uno
switch piu netto verso i codec neurali, ma la valutazione principale per
il dominio immagine usa SSIMULACRA2 come metrica percettiva.
